\documentclass[aspectratio=169]{beamer}
\setlength{\parskip}{5pt}
\usetheme{AnnArbor}
\usecolortheme{default}
\title[ECO247 Presentation]{\small  Data Science, Job Opportunities, Learning R and Data Analysis with R}
\subtitle{\scriptsize Applications in Economics}
\author[Group Seven]{\tiny Shanta \tiny 2019231014 \and
	Torsa \tiny 2019231034 \and
	Sanjida \tiny 2019231051 \and
	Anika \tiny 2019231053 \and
	Apu \tiny 2019231057\and
	Iffath-ara \tiny 2019231063 \and
	Riffat \tiny 2019231075 \\ }
% TODO: \usepackage{graphicx} required
\begin{center}
	\includegraphics[scale=.02]{"../../../Downloads/SUST logo"}
\end{center}
\setbeamercolor{structure}{fg= red}
\setbeamercolor{palette tertiary}{bg=violet, fg=yellow}
\setbeamercolor{palette secondary}{bg=blue, fg=green}
\setbeamercolor{palette primary}{bg=pink, fg=blue}


\setbeamercolor{frametitle}{bg=yellow, fg=blue}

\setbeamercolor{background canvas}{bg=pink!20}


\begin{document}
	\begin{frame}
		\maketitle
	\end{frame}
	\begin{frame}{Objectives of the Presentation}
	\begin{itemize}
		\item  Learn about R-studio make and analysis with data sheet.
		\item Making collaboration R studio with economics.
		\item  Learning about the Job Opportunities.
		\item Learn about R-studio make and analysis with data sheet.
		\item Data security and privacy.
	\end{itemize}
\end{frame}
\begin{frame}{Significance of the Presentation}
	% TODO: \usepackage{graphicx} required
	\begin{center}
		\includegraphics[scale=.4]
		{../../../Downloads/sig}
	\end{center}
	
	
\end{frame}
\begin{frame}{Introduction To Data Science}
	\begin{enumerate}
		\item Data is a collection of facts, like numbers, words, measurements and observations .
		% TODO: \usepackage{graphicx} required
		\includegraphics[scale=.05]{../../../Downloads/xzk3rl06exk-data-types}
		% TODO: \usepackage{graphicx} required
		% TODO: \usepackage{graphicx} required
		\includegraphics[scale=.09]{../../../Downloads/ds-introduction}
		\item As an economic analyst I can use data science to:
		1) Deal with large set 
		2) Summarize data and take economic decisions
		3) Represent economic models easily
		
		
	\end{enumerate}
\end{frame}
\begin{frame}{Importance and Popularity of R and R-Studio and How to Master a Programming Language}
	\begin{enumerate}
		
		\item R and R studio is a free software programming language
		
		\item  Open source allows them to test, experiment, and iterate more than most proprietary software. 
		
		\item  Open source data science tools are most popular because of it is easy to collaborate projects and modify programs in those tools.
		
		\item  R and R studio is a constantly being updated and also growing programming language.
		\item Pick a programming language as per your demand.
		\item Strengthen your fundamentals.
		\item Proceed from basics to advanced level.
		\item 	There is no alternative to practice.
		\item 	Developing  projects and solving problems.
		\item Sharing knowledge and code with others.
		
		
	\end{enumerate}
	
	
\end{frame}
\begin{frame}{Job Opportunities for Economists in the Field of Data Science and Exploring Data Science Jobs}
	\begin{tabular}{|l|l|l|l|l|}
		\hline
		No. &Position&location&Company Name& Salary\\
		\hline
		1 & Data Scientist&Washington, DC
		& Radiant.digital& \$120K- \$125K\\
		\hline
		2&Data Science Director&Newark,NJ&Prudential&\$107K- \$153K \\
		\hline
		3 & Economic Data Analyst&Austin,TX& Realtor.com Careers&\$59K - \$84K\\
		\hline
		4&Senior Data Analyst&Dhaka,BD&Deriv&\$30K-\$45K\\
		\hline
		5&Financial Analyst&Bangladesh&CrossFund&\$35K-\$50K\\
		\hline
		
	\end{tabular}
	\begin{enumerate}
		\item     Create a Compelling Resume and a portfolio.
		\item Create Industry-Specific Projects.
		\item Approach Employers and do networking.
		\item Look For Entry-Level Data Science, Research Intern Jobs.
		\item Consider Working Remotely.
		\item Build Your Personal Brand.
	\end{enumerate}
\end{frame}
\begin{frame}{Tools, Resources and Creating Content for Community and Their privacy and Security}
	\begin{enumerate}
		\item Excel, SPSS, SQL, Python,R,Java are some essential tools to a career in data science.
		\item Youtube channels: Alex the Analyst, Luke Barousse, Ken Jee, Recall by Dataiku, Tina Huang, freeCodeCamp.org, Guy in a Cube.
		\item Analytic Blogs: KDnuggets, Datafloq, Dataeconomy, Data and information visualization, Insidebigdata.
		
		
		\item Content creator can provide their contents to the Data science community by using Github, KDnuggets, Kaggle, LinkedIn, Medium, etc.
		
		
		\item Data security is the concept of protecting digital data from theft, corruption, or unauthorized access.
		
		\item  Data privacy enables individuals to decide and limit access to the use and sharing of their personal data.
	\end{enumerate}
\end{frame}
\begin{frame}{R-Studio Clouding System}
	\begin{enumerate}
		\item R-Studio Cloud is basically a cloud based solution system.
		
		\item R-Studio Cloud runs on a remote Server.
		
		\item R-Studio Cloud allows anyone online to access R-Studio through a web browser.
		
		\item  R-Studio Cloud made it easy to share, teach and learn data science.
	\end{enumerate}
\end{frame}
\begin{frame}{ Designing Online Data Science Training for the Modern Age
	}
	\begin{enumerate}
		\item In modern age, we must go beyond typical e-learning method.
		
		\item The e-learning process should be Outcome-based, comprehensive,iterative and inclusive
		\item Modern learners must have a network for proper guideline online.
		
		
		\item One platform for implementing this method is R-Studio Academy.
		
		
	\end{enumerate}
\end{frame}
\begin{frame}{Necessary Rules for Learning R by Yourself}
	
	\begin{columns} 
		\column{0.5\linewidth} {\color{red!200}Taking Time to Read Books}
		\begin{itemize}
			\item Free ebooks in online
			\item  Code examples, Exercises
			\item Solutions available in online
			\item Include advanced statistical skills and core programming of R.
			\item Give knowledge about web-based data visualization techniques.
		\end{itemize}
		
		\column{0.5\linewidth} {\color{red!200}Using Free Resources}
		\begin{itemize}
			\item Low cost of learning
			\item  R-Studio provide cheat sheets
			\item Free courses available on sites like Coursera, eDX.
			\item Blog post, workshop and resources offered by R groups.
			\item Many institution offer free tutorials, workshops and resources.
		\end{itemize}	
		
	\end{columns}	
\end{frame}
\begin{frame}{Necessary Rules for Learning R by Yourself}
	\begin{columns}
		\column{0.5\linewidth} {\color{red!200}Adopting Good Practices and Being Consistent}
		\begin{itemize}
			\item Consider mindset while learning R
			\item Pseudocode might help at the very beginning
			\item One may choose learning by practically doing code
			\item Commenting is necessary to understand the code in the future
			\item A common style would be the best choice.
			
		\end{itemize}
		
		\column{0.5\linewidth} {\color{red!200}Joining the R Community}
		\begin{itemize}
			\item prolific, with conferences, meetups, and regular online events
			\item   Local UseR or RLadies groups, near an R group location, the R4DS Slack channel
			\item  R bloggers  a great way to stay up-to-date with R n ews.
			\item Online a great way to learn about new features, stay informed,come across tips and tricks and stay motivated.
		\end{itemize}	
		
	\end{columns}		
	
\end{frame}
\begin{frame}{Necessary Rules for Learning R by Yourself}
	\begin{columns}
		\column{0.5\linewidth} {\color{red!200}Reading and Sharing One-another's Code}
		\begin{itemize}
			\item If we find any R product in online, we should run the code to understand the exact benefaction.
			\item Sharing our R code in online will help us become a more proficient coder.
		\end{itemize}
		
		\column{0.5\linewidth} {\color{red!200}Never Box Yourself In}
		\begin{itemize}
			\item Proficiency in R will help us to gain skills that can be applied to other programs.
			\item  The determination we build while learning R can be the begining of venture of our coding journey.
		\end{itemize}	
		
	\end{columns}		
	
\end{frame}
\begin{frame}{Data Analysis Part-1}
	\textbf{1. Smoke Dataset} 
	
	\textbf{2. Ten Variables:} educ, cigpric, white, age, income, cigs, restaurn, lincome, agesq, lcigpric  
	
	\textbf{3. Summary Statistics:}\\
	\textsf{Lincome:}
	\begin{tabular}{ | m{1.5cm} | m{1.5cm}| m{1.5cm} | m{1.5cm} |m{1.5cm} |m{1.5cm} | } 
		\hline
		Min. &  1st Qu. & Median & Mean & 3rd Qu. & Max. \\ 
		\hline
		6.215 & 9.433 & 9.903 & 9.687 & 10.309 & 10.309 \\
		\hline
	\end{tabular}\\
	\textsf{Cigpric:}
	\begin{tabular}{ | m{1.5cm} | m{1.5cm}| m{1.5cm} | m{1.5cm} |m{1.5cm} |m{1.5cm} | } 
		\hline
		Min. & 1st Qu. & Median & Mean & 3rd Qu. & Max. \\
		\hline 
		44.00 & 58.14 & 61.05 & 60.30 & 63.18 & 70.13 \\
		\hline
	\end{tabular}\\
	\textsf{Cigs:}
	\begin{tabular}{ | m{1.5cm} | m{1.5cm}| m{1.5cm} | m{1.5cm} |m{1.5cm} |m{1.5cm} | } 
		\hline
		Min. & 1st Qu. & Median & Mean & 3rd Qu. & Max. \\
		\hline 
		0.000 & 0.000 & 0.000 & 8.697 & 20.000 & 80.000 \\
		\hline
	\end{tabular}
\end{frame}
\begin{frame}{Data Analysis Part-2}
	
	\begin{center}
		ANOVA Hypothesis:
		\[ H_0: Cigpric \ we\ have\ no \ significant\ effect\ on\ lincome \]
		\[H_0: Cigs\ will\ have\ no\ significant\ effect\ on\ lincome\]
		\[H_0: Cigpric\ and\ Cigs\ will\ have\ no\ significant\ effect\ on\ lincome\] 
	\end{center}
	
\end{frame}
\begin{frame}{Data Analysis Part-3}
	
	\begin{enumerate}[I]
		\item Smallest, largest and average years of schooling are 6, 18 and 12.47 respectively 
		\item The number of people don't smoke and claim to smoke 20 cigarettes smoked per day are 497 and 101 respectively. 21.15\% people claim to smoke 20 cigarettes a day 
		\item $$ Shapiro\ Wilk\ Normality\ test\ :\ w\ =\ 0.68741,\ p-value\ <\ 2.2e^{-16} $$
		\item Average income = 19304.83 \\
		The average is not a good measure because there are outliers or extreme value present in the data. \\ 
		Average number of cigarettes smoked per day where restaurant smoking apply if 6.6. 
	\end{enumerate}
	
\end{frame}
\begin{frame}{Data Analysis Part-4}
	\begin{tabular}{ |l |l |l |l |l |l |l |} 
		\hline
		Variable & DF & SS & MS & F-value & P-value \\ 
		\hline
		Cigpric & 1 & 2.3 & 2.3016 & 4.561 & 0.0330 \\ 
		\hline
		Cigs & 1 & 1.8 & 1.8295 & 3.625 & 0.0573 \\ 
		\hline
		Cigpric:Cigs & 1 & 0.0 & 0.0458 & 0.091 & 0.7632 \\
		\hline
		Residuals & 803 & 405.2 & 0.5046 & - & - \\
	\end{tabular}
	
	
\end{frame}
\begin{frame}{Data Analysis Part-5}
	\begin{tabular}{ |l |l |l |l |} 
		\hline
		Variable & P-Value & Decision \\ 
		\hline
		Cigpric & 0.0331 & Significant \\ 
		\hline
		Cigs & 0.565 & Not Significant \\ 
		\hline
		Cigpric:Cigs & 0.7633 & Not Significant \\
		\hline
	\end{tabular}
	
	
\end{frame}
\begingroup
\setbeamertemplate{footline}{}
\begin{frame} [c] { }
	
	\centering
	\large  Thank You! \\
	\centering
	\bigskip
	Questions and Answers?	
\end{frame}
\endgroup

\end{document}